% ===========================================================================
% Section 2: Background and Related Work (~1.5 pages)
% ===========================================================================

% --- 2.1 Darshan I/O Instrumentation ---
% \subsection{Darshan I/O Profiling}
% \label{sec:related:darshan}
%
% TODO: Darshan overview, modules (POSIX, MPI-IO, STDIO), counters
% Cite: \cite{snyder2016darshan}
% Polaris dataset: \cite{snyder2025darshan}

% --- 2.2 I/O Characterization Studies ---
% \subsection{HPC I/O Characterization}
% \label{sec:related:characterization}
%
% TODO: Prior I/O characterization work
% Cite: \cite{luu2015multiplatform, paul2022characterizing, bang2020clustering, bez2023survey}

% --- 2.3 ML-Based I/O Analysis ---
% \subsection{ML-Based I/O Analysis}
% \label{sec:related:ml}
%
% TODO: AIIO, WisIO, Kim et al.
% Cite: \cite{dong2023aiio, devarajan2025wisio, kim2023predicting}

% --- 2.4 LLM-Based Approaches ---
% \subsection{LLM-Based I/O Diagnosis}
% \label{sec:related:llm}
%
% TODO: IOAgent, ION
% Cite: \cite{zheng2025ioagent, lockwood2024ion}

% --- 2.5 ML+LLM Integration in Systems ---
% \subsection{ML+LLM Integration}
% \label{sec:related:hybrid}
%
% TODO: RCACopilot and other ML+LLM patterns
% Cite: \cite{chen2023rcacopilot}
% Explain why this pattern is emerging and how we adapt it for I/O
